\section{LlamaIndex: construcción de un asistente de conocimiento multimodal}

\subsection{De Basic RAG a Advanced RAG}

Jerry Liu señala cuatro limitaciones fundamentales del Basic RAG:
\begin{enumerate}
    \item \textbf{Procesamiento de datos tosco}: la simple división del texto en chunks no puede manejar tablas, imágenes ni diseños complejos
    \item \textbf{El LLM se usa solo para la síntesis}: no se aprovechan las capacidades de razonamiento y planeación
    \item \textbf{Sin estado}: cada consulta es independiente, sin memoria personalizada
    \item \textbf{Salida única}: solo puede generar respuestas de texto simples
\end{enumerate}

\begin{warningbox}{Garbage In Garbage Out}
La calidad de una aplicación de RAG depende de la calidad de su pipeline de procesamiento de datos. Si el analizador de PDF no puede extraer correctamente las tablas y las imágenes, ningún razonamiento posterior del LLM, por potente que sea, puede compensarlo. La calidad de los datos es una condición necesaria para las aplicaciones de LLM en producción.
\end{warningbox}

\subsection{Pipeline de RAG multimodal}

Pasos clave para construir un RAG multimodal:

\begin{enumerate}
    \item \textbf{Análisis de documentos de alta calidad}: extraer por separado el texto, las tablas, las gráficas y las imágenes de un PDF como elementos estructurados
    \item \textbf{Indexación jerárquica}: generar para cada elemento múltiples representaciones textuales (resúmenes, sub-chunks, etc.) que apunten al elemento fuente
    \item \textbf{Recuperación multimodal}: dada una consulta, recuperar chunks de texto e imagen, aprovechando un LLM multimodal para procesar simultáneamente la información textual y la visual
\end{enumerate}

\begin{knowledgebox}{Nuevas oportunidades de los modelos multimodales}
Los modelos dominantes actuales (GPT-4o, Claude 3.5 Sonnet, Gemini) admiten la entrada simultánea de texto e imágenes. Esto significa que, incluso si la indexación se construye con un modelo de embeddings de texto, en la etapa final de inferencia aún es posible alimentar la imagen original al LLM y lograr una comprensión visual genuina.
\end{knowledgebox}

\subsection{Agentic RAG: agregar una capa de razonamiento sobre el RAG}

\begin{importantbox}{La idea central del Agentic RAG}
Tratar la recuperación de RAG como una herramienta (tool) y agregar por encima una capa de razonamiento de Agent. El Agent puede decidir cómo descomponer el problema, qué herramientas de recuperación invocar y cuándo hacer una reflexión de verificación, de modo que pueda abordar problemas complejos que el RAG básico no puede responder (resúmenes, comparación de múltiples documentos, análisis de varios pasos, etc.).
\end{importantbox}

Dos opciones de arquitectura:
\begin{itemize}
    \item \textbf{Flujo restringido}: el desarrollador define manualmente la lógica if-else y las reglas de enrutamiento; es más confiable pero menos flexible
    \item \textbf{Arquitectura de Agent general} (como ReAct): el Agent elige por sí mismo las herramientas y la ruta de razonamiento; es más flexible pero menos confiable y de mayor costo
\end{itemize}

Regla práctica: en una arquitectura de Agent general, el número de herramientas debe mantenerse en lo posible entre 4 y 5, sin superar las 10.

\subsection{Despliegue en producción}

LlamaIndex ofrece un sistema de orquestación basado en eventos que permite:
\begin{itemize}
    \item Encapsular cada flujo de trabajo de Agent como una API de microservicio
    \item Comunicación entre Agents mediante una cola de mensajes central
    \item Human-in-the-loop: el Agent puede pausarse para esperar la entrada de un humano antes de continuar la ejecución
    \item Soporte para desarrollo local y despliegue en Kubernetes
\end{itemize}

\subsection{Resumen de la sección}

LlamaIndex, partiendo de la calidad del procesamiento de datos, construye un pipeline completo que va del análisis de documentos a la recuperación multimodal y de ahí al razonamiento Agentic. Su idea central es: una buena calidad de datos es la base, el razonamiento del Agent es la superestructura, y ninguna de las dos puede faltar.

